\newpage
\section{Introduction}

Plasma physics research has been strongly motivated by controlled fusion for more than 70 years, with contemporary developments focusing on magnetic confinement (tokamaks, stellarators) and inertial confinement (lasers, Z-pinches). NERS research emphasizes intense laser–plasma interactions (ZEUS), pulsed-power and high-power microwave sources, as well as low-temperature plasmas (GEC).

\section{What is a plasma?}
A plasma is an ionized gas and is sometimes referred to as the "fourth state of matter". This course emphasizes high-temperature plasmas (fusion plasmas, > 1 keV) rather than low-temperature plasmas (processing plasmas, $\sim$ 1 eV).

Three fundamental parameters characterize a plasma:
\begin{itemize}
  \item the particle density $n$ (measured in particles per cubic meter),
  \item the temperature $T$ of each species (usually measured in eV, where 1 eV = 11605 K),
  \item the steady-state magnetic field $B$ (measured in Tesla).
\end{itemize}

A host of other properties can be derived from these three fundamental parameters.

\subsection{Length Scales}

The length scales we'll be dealing with in this course are:
\begin{itemize}
  \item $a_0$ = atomic scale ($\sim = 10^{-8}$ cm $= 10^{-10}$ m)
  \item $n$ = charged particle number density ($1/\text{cm}^{-3}$ (interstellar medium) $\longrightarrow 10^{20}/\text{cm}^{-3}$ (Z-pinch, laser fusion, laser-wakefield plasmas))
    \begin{itemize}
      \item Most laboratory plasmas have charged particle densities on the order of 
    \end{itemize}

  \item $\frac{1}{n}$ = the average volume occupied by a single particle
  \item $n^{-\frac{1}{3}}$ = average interparticle separation
\end{itemize}

In a gaseous state, the average interparticle distance should be much greater than the Bohr radius ($n^{-\frac{1}{3}} \gg a_0$. This is because atoms and ions will behave mostly independently if further apart than the Bohr radius, which dictates the natural scale of bound electron orbitals. In an ionized state, the thermal energy must be greater than or equal to the ionization energy ($k_BT \geq E_i$).

\begin{prop}
  A plasma necessarily has a large number of electrons within a Debye sphere.
\end{prop}

\begin{eg}
  Tokamak Plasma
\end{eg}
\begin{explanation}
  In these plasmas, 
\end{explanation}

\subsection{Other Properties of a Plasma}
\begin{enumerate}
  \item Particle kinetic energy is \textit{much} greater than the average coulombic interaction energy.
  \item Over a length scale $L \gg \lambda_D$, plasma is neutral.
  \item Particles are scattered mostly by multiple, cumulative small angle collisions rather than by large angle collisions.
  \item Collective interactions dominate over coulombic interactions between individual particles.
  \item One may ignore the discreteness of individual particles and can treat plasma motion with a continuum description (e.g. Vlasov, 2-fluid, MHD)
\end{enumerate}

\begin{definition}
  Plasma dynamics is determined by the \textit{self-consistent interaction between electromagnetic fields and statistically large numbers of charged particles}.
\end{definition}
